\section{Related Work}
\label{sec:related_work}

Document-level relation extraction (DocRE) requires reasoning across sentence boundaries to identify relations between entity pairs.
DocRED \citep{yao2019docred} established the first large-scale benchmark for this task, catalyzing a line of encoder-based methods that treat DocRE as classification over entity pair representations.
ATLOP \citep{zhou2021atlop} introduced adaptive thresholding for multi-label prediction and localized context pooling for entity representations.
A parallel thread incorporates sentence-level evidence into these architectures.
SAIS \citep{xiao2022sais} introduced evidence supervision as an auxiliary training signal and showed that joint evidence prediction improves extraction through data augmentation and ensemble inference.
DREEAM \citep{ma2023dreeam} extended this approach with dynamic memory-based evidence retrieval that guides attention to supporting sentences.
EIDER \citep{xie2022eider} jointly trains an evidence extractor and fuses evidence representations into the relation classifier.
In all three systems, evidence identification is binary sentence classification over a fixed candidate set, a formulation that adds no output complexity.
Recent work has recast DocRE for large language models.
AutoRE \citep{autore2024} fine-tunes LLMs with QLoRA using a relation-head-facts decomposition; \citet{xu2024lmrc} frame extraction as multi-choice reading comprehension; RelPrior \citep{wang2025relprior} uses relation priors to guide LLM-based extraction; GLiDRE \citep{zhang2025glidre} presents a generalist lightweight approach; Ru-LLM \citep{wang2025rullm} incorporates logical rule constraints; and \citet{chen2025docksrag} combine retrieval-augmented generation with structured knowledge.
These methods match or surpass encoder baselines but produce bare triplets without evidence. Among LLM-based generative methods at the 8B scale, our RSFT model (54.33 Ign-F1 on DocRED dev) outperforms AutoRE (52.14\footnote{AutoRE Ign-F1 is from the LLaMA3-8B reproduction by \citet{wang2025relprior}; the original paper reports only F1.}) and is slightly below RelPrior ($-$1.17pp), while being the only method that additionally provides sentence-level evidence grounding.
The critical distinction lies in computational nature: in SAIS and DREEAM, evidence retrieval is classification that regularizes the encoder without increasing output length; in LLM-based DocRE, evidence requires additional output tokens per triple, imposing costs that depend on model capacity and optimization strategy.
Whether this cost is inevitable or can be mitigated is the central question of our work.

Group Relative Policy Optimization (GRPO; \citealt{shao2024deepseekmath}) eliminates the separate value network by computing advantages within groups.
Subsequent work addresses failure modes: DAPO \citep{yu2025dapo} prevents entropy collapse, Dr.\ GRPO \citep{liu2025drgrpo} corrects length bias, and multi-objective variants \citep{liu2026gdpo,ito2025mogrpo} handle competing reward signals.
Applications to structured IE include R1-RE \citep{dai2025r1re} and REAR \citep{wang2025rear}.
Direct Preference Optimization \citep{rafailov2023direct} reformulates RLHF as implicit reward maximization.
Our decoupling spectrum is distinct from catastrophic forgetting in RL alignment \citep{jin2025rl_panacea}, where models lose pre-trained capabilities across tasks; in our setting, extraction and evidence compete \emph{within} a single task, resembling negative transfer in multi-task learning \citep{ruder2017overview,standley2020tasks,yu2020gradient}, amplified by QLoRA's constrained parameter space \citep{gao2023scaling,moskovitz2024constrained}.
RSFT avoids this intra-task competition by operating on discrete outputs rather than gradients, extending the sample-and-select principle \citep{wang2023selfconsistency} to structured IE with joint quality scoring.

Evidence faithfulness evaluation has received increasing attention: \citet{laban2022summac} re-visit NLI-based models for summarization inconsistency detection, \citet{li2025drift} propose dual-reward inference for faithful rationale generation, \citet{li2025srvf} use supervised rationale verification with contrastive feedback for relation extraction, and \citet{zhang2024infuse} introduce fine-grained NLI-based faithfulness metrics for summarization.
Our CED scoring builds on similar NLI-based principles but targets evidence-relation entailment rather than free-form rationales.

\citet{zelikman2022star} introduced STaR for iterative self-improvement; \citet{dong2023raft} generalize this with reward-model-based sample filtering; recent analysis shows rejection sampling implicitly performs PPO-like optimization \citep{tang2025rsftppo}.
Our RSFT is a single-round variant applied to evidence-grounded DocRE with CED/NLI scoring as a domain-specific quality filter.
Because RSFT selects complete outputs rather than differentiating through a reward signal, it cannot exploit the extraction-evidence trade-off that gradient methods discover, inherently preserving joint quality.
